\section{Objetivos}
\subsection{Objetivo general}
Comparar experimentalmente ANIA y L3C-DeepMFC para la cancelación de realimentación acústica en tiempo real, considerando ganancia adicional estable, calidad de voz, robustez y costo computacional bajo condiciones de latencia controlada.

\subsection{Objetivos específicos}
\begin{enumerate}
\item Definir un modelo del lazo acústico y un protocolo reproducible que especifique señales, caminos, ganancias, frecuencias de muestreo y condiciones de cambio del entorno.
\item Implementar y verificar la identificación y cancelación de ANIA, distinguiendo la calibración manual del control automático descrito en la publicación original.
\item Implementar, entrenar y verificar L3C-DeepMFC a partir de su descripción publicada, con separación entre entrenamiento, validación y evaluación externa.
\item Integrar ambos métodos en un banco común de procesamiento continuo y comprobar la equivalencia entre las versiones de referencia y las versiones utilizadas en tiempo real.
\item Medir estabilidad, calidad de voz y recuperación ante cambios del camino mediante ensayos emparejados, con incertidumbre y condiciones de validez explícitas.
\item Cuantificar latencia, tiempos de procesamiento y uso de recursos en PC, y evaluar la factibilidad de ejecución sostenida en Raspberry Pi 4.
\end{enumerate}

\section{Justificación}
La contribución esperada será una comparación reproducible que vincule el desempeño acústico con las restricciones de implementación. La ausencia de acople, por sí sola, no acreditará una solución satisfactoria: un método podría evitarlo atenuando excesivamente la señal útil. De manera similar, una mejora de calidad fuera de línea podría perder utilidad si requiere un retardo excesivo o incumple los plazos de procesamiento.

El uso de un banco común permitirá documentar estos compromisos y separar errores del algoritmo de limitaciones del dispositivo o del controlador de audio. La verificación numérica previa reducirá el riesgo de atribuir a un método diferencias causadas por una implementación incorrecta. La evaluación externa con grabaciones propias permitirá examinar condiciones distintas de las utilizadas para ajustar parámetros.

El trabajo aportará implementaciones documentadas, un protocolo de medición y evidencia que permita seleccionar un método según la aplicación. El aporte seguirá siendo válido si el modelo neuronal no supera al método clásico o si no resulta viable en la plataforma embebida, siempre que las conclusiones se sustenten en mediciones comparables.
